\paragraph{Conjugate-coordinate class times a fuzzy factor: response of the scale of \(K\) to the frame, the exponent of the \(K\) masses and the scalar channel of the static response (conditional; bounded obstruction).}
This paragraph supplements the preceding one, in which the coupling of massive modes to the scale of \(K\) was recorded as the two alternatives (S1), (S2) and not derived. Keep its setting and notation, with \(V=T\), equal weights and \(q\cdot y=0\) where they are used, and write \(X^i=\chi\Lambda^i\) with fixed Hermitian \(\Lambda^i\), \(\bar a=-\tfrac14\operatorname{tr}[\Lambda^i,\Lambda^j]^2\ge0\), \(\bar n=\operatorname{tr}\Lambda^i\Lambda_i\), \([\Lambda^i,[\Lambda_i,\Lambda_A]]=\mu_A^2\Lambda_A\), \(\gamma^{\mu\nu}=\eta_{ab}E^{a\mu}E^{b\nu}\), \(g=\gamma^{-1}\) and \(\phi=\tfrac12\ln(\rho_M/\sqrt{|g|})\), with \(\rho_M\) taken proportional to the density \(\rho\) of the preceding paragraphs, so that \(\phi\) is the field there up to a constant. A mode \(\psi\Lambda_A\) has the mass \(M^2=\chi^2\mu_A^2\) in the operator \(\gamma^{\mu\nu}\partial_\mu\partial_\nu\) of the preceding paragraph. Define the exponent \(\beta_K\) by \(\delta\ln M^2=2\beta_K\delta\phi=\beta_Kt/2m^2\). On the mass shell the source of a free massive mode is then \(X_\beta=T^{\mathrm m}-\tfrac12\beta_K\eta\operatorname{tr}T^{\mathrm m}\), which couples to \(h_\beta=\bar h-\tfrac12\beta_K\eta t\), the perturbation of \(e^{-2\beta_K\phi}\gamma\); (S1) is \(\beta_K=1\) and (S2) is \(\beta_K=0\).

\emph{Tree level.} For a product the operator on \(\psi\Lambda_A\) is the sum of the operator on \(M\) and \(\chi^2\mu_A^2\), and, with \(g^2\) in \(\rho_M/g^2\) the coupling of the tree form of the preceding paragraph, the tree form is exactly \(\tfrac{\rho_M}{2g^2}\operatorname{tr}(\Lambda_A^2)(\gamma^{\mu\nu}\partial_\mu\psi\partial_\nu\psi+\chi^2\mu_A^2\psi^2)\). So at fixed \(\Lambda^i\) and \(\chi\) the mass is rigid in \(\gamma\), which is (S2), and in the frame \(e^{-2\phi}\gamma\), the effective metric of~\cite{Steinacker2026}, it is \(e^{-2\phi}\chi^2\mu_A^2\), the factor stated in~\cite{MantaSteinacker2025b}. At fixed \(\chi\) a loop carries by power counting one factor \(g^2\sqrt{|g|}/\rho_M=g^2e^{-2\phi}\), a trace over \(M\) against the weight \(\rho_M/g^2\), so corrections to \(M^2\) at \(\ell\) loops have the exponent \(-\ell\). So, while the loop expansion is controlled, a value \(\beta_K\) of order one can come only from the response of \(\chi\); a loop correction of negative sign contributes a small positive part.

\emph{Scaling.} Under \(X^a\to\sqrt\lambda X^a\), the operators of the content of the preceding paragraphs are exactly \(\lambda H_M+\chi^2H_K\). Since the supertrace of the identity vanishes, the one-loop functional, where its proper-time form converges, depends on \(\lambda,\chi\) only through \(\chi^2/\lambda\) (random product matrices, agreement to twelve digits; the invariance under a common rescaling is used in~\cite{MantaSteinacker2025b}). At one loop alone \(\delta\ln\chi=\tfrac12\delta\ln\lambda\) is a flat direction, and for this homogeneous mode \(\phi=\ln\lambda\) up to a constant, so its exponent is \(\tfrac12\).

\emph{Potential of \(\chi\).} From the factorisation of the preceding paragraphs, with equal weights, \(F=(n+1)^2\) and hats for the unit scale \(\chi=1\), the terms of the action density without derivatives of the wave that depend on \(\chi\), through first order in the background torsion, are
\[
\frac{\rho_M}{g^2}V_{\mathrm t}(\chi)+U_4+U_2,\qquad
U_4=-\frac{F}{32\pi^2}\sqrt{|g|}\,\hat U\chi^4,\qquad
U_2=-\frac{F}{32\pi^2}\sqrt{|g|}\Bigl[2(\hat V-\hat T)-\tfrac34\vartheta(\hat V-2\hat T)\Bigr]\chi^2,
\]
with \(\hat U=\int_0^\infty\mathrm ds\,s^{-3}(\hat v^K-\hat\tau^K)(s)\), and \(V=F\hat V\chi^2\), \(T=F\hat T\chi^2\) at equal weights. Here \(V_{\mathrm t}\) contains \(\bar a\chi^4\) and any supplied stabiliser, and \(U_2\) uses \(R/6+E_0=2-\tfrac34\vartheta\) and the index density \(-6\vartheta\) of the background. This has the structure of Equation~(7.48) of~\cite{Steinacker2023}, including the negative sign of the torsion term in the Lorentzian section. At the next order the first term \(s_0^K(v^M-2\sigma^M)\) of the factorisation enters: with the fibre flux dropped, at the frozen symbol \(v^M-2\sigma^M\) has the value \(\tfrac{15}{4}\vartheta^2F\sqrt{|g|}/16\pi^2\) at order \(s^0\) on the background, an order that has no term linear or quadratic in the curvature, as recorded in the preceding paragraphs. So the derivative of the density with respect to \(\ln\chi\) gains \(Y=\tfrac{15F}{64\pi^2}\vartheta^2\sqrt{|g|}N_K\ge0\), with \(N_K\) the number of modes of \(K\) with \(\chi^2\mu_A^2\gg\vartheta\). It is not suppressed by the flux of \(K\) or at large \(N\); on \(S^2_N\) at large \(N\), \(Y/2|U_2|=\tfrac54\vartheta/\chi^2\). The fibre flux in \(Y\) is treated below. The background Feynman gauge is used (Hessian checked by finite differences); its dependence at the stationary point is of higher order by the usual argument, not computed. The value \(U_2\) above is the order \(\vartheta\) of the mixed term, and on \(S^2_N\) the mixed term with all orders of the frozen symbol agrees with it (ratio above \(0.9999\) at \(N\ge256\) for \(\chi^2\ge0.1\vartheta\)). At \(q\cdot y=q\cdot v=0\) the product \(\bar T\cdot\delta T\) vanishes, so with \(\bar h^{22}=\bar h^{11}\) one has \(\delta\ln(\sqrt{|g|}X)=\Omega/m^2\), \(\Omega=\tfrac12\bar h^{00}-\tfrac53\bar h^{11}\), and \(\Omega=-\tfrac12t\) where \(\bar h^{11}\) is negligible, which holds at large \(\tau\) (below). Write \(U\) for this density with the first term and \(U_{\chi\chi}=\partial^2U/\partial(\ln\chi)^2\) at a stationary point. Since \(\bar T^{ab\mu}y_\mu=0\), a perturbation along \(y^\mu y^\nu\) leaves the Gaussian average over the loop momentum unchanged, so where \(\bar h^{11}\) is negligible \(U_4\), \(U_2\) and the fibre-constant part of \(Y\) carry \(\sqrt{|g|}\) alone. For \(q^2\) below the mass of \(\chi\) defined below and \(V=T\), the first-order response gives in this static channel at large \(\tau\)
\[
\beta_K=\frac{2(4U_4+2U_2+Y)}{U_{\chi\chi}}=-\frac{2\,\partial_{\ln\chi}(\rho_MV_{\mathrm t}/g^2)}{U_{\chi\chi}},
\]
and for the homogeneous mode, with \(\partial Y/\partial\ln\chi\) neglected, \(\beta_K=(8U_4+2U_2)/U_{\chi\chi}\), which is \(\tfrac12\) for \(V_{\mathrm t}=0\). Hence at a local minimum, \(U_{\chi\chi}>0\), \(\beta_K\) has the sign opposite to the slope of the tree term; for \(U_4\le0\) and \(U_2<0\) one has \(\beta_K<0\) for every \(V_{\mathrm t}\) exactly where \(Y<4|U_4|+2|U_2|\), in particular for \(\chi^2\hat T\gg\vartheta N_K\). With \(\tilde U=\rho_M\bar a\chi^4/g^2\) and \(\varepsilon=U_4/\tilde U\): for \(V_{\mathrm t}=\bar a\chi^4\), \(\beta_K=-1/(1+\varepsilon-Y/4\tilde U)\); adding a tree cubic term, of the form of the cubic term of Equation~\eqref{eq:Smatrix} read on \(K\) and supplied there, gives \(\beta_K=2(4U_4+2U_2+Y)/(4\tilde U+4U_4-2U_2-3Y)\); adding a tree mass term gives \(\beta_K=(4U_4+2U_2+Y)/(4\tilde U+4U_4-Y)\), with \(\partial Y/\partial\ln\chi\) neglected in these three forms. Where the tree stabiliser dominates, \(\beta_K\to0\), from below where the slope of the tree term is positive and from above where it is negative. For \(Y\) negligible, \(\beta_K>0\) needs \(U_4>0\), so \(\hat U<0\); then \(\beta_K=1\) is attained in the cubic case at \(U_4=\tilde U-\tfrac32U_2\) and approached from below in the mass case as \(U_4\) comes to dominate \(\tilde U\) and \(U_2\), each with a tree term of negative slope. In the rotating class of~\cite{MantaSteinacker2025b}, Equation~(3.37) there gives \(\chi\propto\lambda^{-1/6}\) at fixed charge and fixed \(\tau\), the exponent \(-\tfrac16\) for the homogeneous mode; the fluctuation operator about a rotating \(K\) is not treated here.

\emph{Fibre flux in the first term.} On the background \(\mathbb F_{ab}\) is the action of the generators \(J_{ab}\) of \(\mathfrak{so}(3,1)\), which commute with \(\Box\). Write \(\mathfrak{so}(3,1)\) over \(\mathbb C\) as \(\mathfrak{su}(2)\oplus\mathfrak{su}(2)\), so that the vector is \((\tfrac12,\tfrac12)\) and the Dirac spinor is \((\tfrac12,0)\oplus(0,\tfrac12)\). Write \(H_i=H_0+Z\) for the vector and for the Dirac index, so that \(Z=C_{31}^{\mathrm{tot},i}-C_{31}-C_{\Sigma_i}\) in the form of the preceding paragraphs, with \(C_{31}=\tfrac12J_{ab}J^{ab}\) for real generators; the sign is fixed by \(H_VT_a=3H_0T_a\), that is \(Z=6\) on the singlet of the orbital \(T_a\) with the vector index. On an irreducible orbital component with labels \(j_\pm\), \(\Xi_\pm=2j_\pm+1\), and per state, with \(P(\Xi)=e^{-s}[2\cosh s\Xi+2\Xi^{-1}\sinh s\Xi]\), exactly,
\[
\operatorname{tr}_De^{-sZ}=P(\Xi_+)+P(\Xi_-),\qquad
\operatorname{tr}_Ve^{-sZ}=P(\Xi_+)P(\Xi_-),\qquad
\operatorname{tr}_Ve^{-sZ}-2\operatorname{tr}_De^{-sZ}+4=\bigl(P(\Xi_+)-2\bigr)\bigl(P(\Xi_-)-2\bigr)
\]
(on each chiral block \((Z-1)^2=\Xi^2\) and the trace of \(Z\) over the block vanishes, in any representation; explicit vector and gamma matrices at finite-dimensional labels). On a unitary principal-series component \(\Xi_\pm=\pm m_L+\mathrm i\varsigma\), with \(m_L\) an integer and \(\varsigma\) real, so the two factors are complex conjugates and the integrand of the first term is \(|P(\Xi_+)-2|^2\ge0\) on each component, for every \(s\) and with every fibre flux included. At fibre order \(J\), at a frozen base point, the rotations about \(y\) act as spin \(J\), \(|m_L|\le J\), the density of \((m_L,\varsigma)\) is the Plancherel density \(\varsigma^2+m_L^2\), the squared boost is \(\varsigma^2+1-m_L^2+c_J\) with \(c_J=J(J+1)\), and \(H_0\) is the squared boost divided by \(\vartheta\) plus \(\mu_J\), with the cross term dropped and the direction along \(y\) a spectator. With \(M_A^2=\chi^2\mu_A^2\), the proper time is small where \(M_A^2+\mu_J\gg\vartheta\), and there the average of the integrand over these labels is \(s^2[\tfrac{15}4\vartheta^2-s\vartheta c_J+s^2c_J^2]\), positive, with corrections of relative order \(s\) and \(s\vartheta\) that do not depend on \(J\) at this order (the integral over \(m_L\) is exact in the limit of large \(J\) at fixed \(sJ^2\); the exact average agrees). Hence
\[
Y=\frac{\sqrt{|g|}}{32\pi^2}\sum_A\sum_{J\le n}(2J+1)\,2M_A^2\Bigl[\frac{15\vartheta^2}{4\mathcal D}-\frac{\vartheta c_J}{\mathcal D^2}+\frac{2c_J^2}{\mathcal D^3}\Bigr],\qquad\mathcal D=M_A^2+\mu_J,
\]
the sum running over the modes with \(M_A\ne0\) and \(\mathcal D\gg\vartheta\), every summand positive since the bracket has negative discriminant in \(\mathcal D\), which reduces to the value above for \(M_A^2\gg\mu_n\). On \(S^2_N\) its relative correction is of order \((n/\chi N)^2\) with a logarithm, the condition for equal weights below (the ratio to the value above is \(0.998\) at \(N=1024\), \(n=10\), \(\vartheta=1\), \(\chi^2=0.3\)), and the fibre-flux part of the mixed term contributes \(4(2J+1)c_J\) at level \(J\) to \(\partial U/\partial\ln\chi\) in the same units for \(\mu_J\gg\chi^2\), and less for smaller \(\mu_J\) (numerical), of relative order at most \((n/N)^2\). The cross term commutes with \(J_{ab}\), so it leaves the non-negativity of the integrand unchanged, but it changes the weights at relative order \(sJ^2\), the order of the \(J\)-dependent terms of this form: those terms hold only with it dropped, and what carries over is their order. So at equal weights the fibre flux does not change the leading value of \(Y\) or its sign on the background. The response of the fibre levels to \(\bar h\) and the part of the first term with two fibre fluxes and two base derivatives of the wave are not computed.

\emph{Explicit \(K\).} For the fuzzy sphere \(S^2_N\), \(\Lambda^i=J^i\) in three of the six directions, the spectra are closed: the vector \(l(l-1)\), \(l(l+1)+2\), \((l+1)(l+2)\) and the spinor \(l^2\), \((l+1)^2\), without negative mode; \(\bar a=N(N^2-1)/8\), \(\bar n=N(N^2-1)/4\), and \(\chi\) is the only singlet of its symmetry group, so the masses share one exponent at this order. Numerically \(\hat U>0\) with \(\hat U/N^2\to1\), \(\hat T/N^2\to4\) and \(\hat V-\hat T\simeq4\ln N+1.016\) (exact spectra, \(N\le512\)), so \(U_4<0\), \(U_2<0\) and, at large \(N\) with \(x=\chi^2/\vartheta\) in this part, not the ratio \(a/b\) of the preceding paragraphs, and \(N^2\chi^2\gg\vartheta\), \(4U_4+2U_2+Y\) is proportional to \(-4x^2-6x+\tfrac{15}2\), positive for \(x<(\sqrt{39}-3)/4\simeq0.81\). For \(V_{\mathrm t}=\bar a\chi^4\) a stationary point then exists only for \(|\varepsilon|\ge\tfrac{10}{13}\) and a local minimum only for \(\tfrac{10}{13}<|\varepsilon|<1\), with \(x>\tfrac52\) and \(\beta_K<-1\), so there is none at weak coupling. With a tree cubic term \(\beta_K=1\) is attained at a local minimum with \(\hat U>0\): in units \(F\sqrt{|g|}N^2\vartheta^2/32\pi^2\), at \(x=\tfrac3{10}\), a tree quartic \(7.935\) and a tree cubic \(-12.36\) give \(U_{\chi\chi}=10.68\) and \(|\varepsilon|\simeq0.011\). These tree values are supplied, and the tree quartic there is of the order of \(Y\), so the one-loop term is not small against the tree term, and higher loops are not computed. The sign \(\hat U>0\) holds also for the sampled \(S^2_N\times S^2_{N'}\) and the squashed \(\mathrm{su}(3)\) orbits (a finite list, no general statement). For modes between stacks with constant relative fluxes \(f_1\ge f_2,f_3\ge0\), exactly, \(v-\tau=16s_0\prod_\pm\sinh\tfrac12(f_1\pm f_2\pm f_3)s\), a product over the four sign choices, so for \(f_1<f_2+f_3\), the range without a negative mode, the integrands of \(\hat U\) and \(\hat V-\hat T\) are negative for every \(s\), and for \(f_1=f_2+f_3\), \(v\equiv\tau\). These integrals diverge at small \(s\) for such non-compact stacks, so \(\hat U<0\) and \(\hat V<\hat T\) are signs with a lower limit on \(s\); the scaling \(U_4\propto\chi^4\) then needs a limit that scales as \(\chi^{-2}\), such as a fixed truncation of the modes, for which the sign is not established. On this evidence the sign of \(\hat U\) depends on \(K\).

\emph{Scale dependence.} With the tree kinetic coefficient \(Z_\chi=\rho_M\chi^2\bar nm^2/g^2\) of \(\ln\chi\) (the loop part is not computed), \(\ln\chi\) has the mass \(m_\chi^2=U_{\chi\chi}/Z_\chi\) in \(q^2\), and integrating it out replaces \(\beta_K\) by \(\beta_Km_\chi^2/(q^2+m_\chi^2)\), with a contact term \(-2\varrho_0^2m^4/Z_\chi(q^2+m_\chi^2)\). So \(M^2\) is rigid in \(\gamma\) above \(m_\chi\), for every \(K\). For \(V_{\mathrm t}=\bar a\chi^4\), \(m^2m_\chi^2=8(1+\varepsilon-Y/4\tilde U)\bar a\chi^2/\bar n\), which is \(4(1+\varepsilon-Y/4\tilde U)\chi^2\) on \(S^2_N\). So (S2) holds above \(m_\chi\) and approximately where a tree stabiliser dominates, and below \(m_\chi\) the value (S1) is a condition on the tree stabiliser and on \(\chi^2/\vartheta\), not an output of this class.

\emph{Scalar channel of the static response.} For dust and \(X_\beta\), with \(g_2=\vartheta^{-2}[1-\tfrac32c_0\operatorname{arctanh}c_0-\tfrac12(\vartheta-1)]\) the finite part of \(\langle(v^0)^4/(v\cdot v)^2\rangle\), in addition to the function \(g=g(\tau)\) of the preceding paragraph (in the next display \(g^2\) is its square, not the coupling),
\[
\frac{I[X_\beta]}{\varrho_0^2}=\varrho\bigl(\tfrac12\beta_K-g\bigr)^2+g-g_2+\varrho\,(g_2-g^2),
\]
with \(g-g_2\simeq2/\sinh^2\tau\) and \(g_2-g^2\simeq-\tfrac43\sinh^{-4}\tau\), which tends to \(\varrho(1-\tfrac12\beta_K)^2\) as \(\tau\to\infty\) (regulated quadrature agrees). So the response relative to the one-coefficient value is \(4g(1-\beta_K)+\beta_K^2\) at \(\varrho=1\), which reproduces \(1\) and \(4g\) for (S1) and (S2). At \(\varrho=1\), \(\bar h^{00}\), \(\bar h^{11}\) and \(t\) are proportional to \((2-\beta_K)g\), \(\tfrac12\beta_K(1-g)\) and \(\beta_K-2g\), and the ratio of the potentials read in the variable to which the source couples is
\[
\frac{h_\beta^{11}}{h_\beta^{00}}=\frac{\beta_K(1+g-\beta_K)}{4g(1-\beta_K)+\beta_K^2}\ \longrightarrow\ \frac{\beta_K}{2-\beta_K},
\]
while \(\bar h^{00}+\bar h^{11}\) does not depend on the reading. At large \(\tau\) the source therefore acts through \(t\) alone: with \(L=32\pi^2m^6/A\) the reduced form is \(S_E=(q^2/4L)t^2-(1-\tfrac12\beta_K)\varrho_0t\), the phase average of the density \((A/4\pi^2)\sqrt{|g|}g^{\mu\nu}\partial_\mu\phi\partial_\nu\phi\) in this static channel, with the positive sign in \(S_E\); \(t\) is gauge invariant here, since \(q\cdot v=0\), and the matter charge is \(1-\tfrac12\beta_K\). With a term \(\tfrac12\mathcal Mt^2\), \(q_*^2=2\mathcal ML\), at \(\varrho=1\) the response is \(-L\varrho_0^2[(\tfrac12\beta_K-g)^2/(q^2+q_*^2)+g(1-g)/q^2]\). A heavy \(t\) therefore leaves only the part with \(g(1-g)\simeq2/\sinh^2\tau\) and ratio \(\tfrac12\) for every \(\beta_K\). A light \(t\) gives \(\beta_K/(2-\beta_K)\), so the value \(1\) at large \(\tau\) needs \(\beta_K=1\), and \(1-\beta_K/(2-\beta_K)=2(1-\beta_K)/(2-\beta_K)\).

\emph{Vacuum term.} At a fibre point the frame perturbation has rank one, so \(\sqrt{|g_n|}=m^{-4}(1+x_n)^{-1}\) with \(x_n\) proportional to \(\alpha\cdot v/m\) exactly, and for a real wave the second-order part of \(U_4\) is \(\tfrac12V_0m^{-6}(\alpha\cdot v)^2\), \(V_0=-F\hat U\chi^4/32\pi^2\), not the weight \(V_0\) of the preceding paragraphs. This is the expansion of Equation~(4.27) of~\cite{Steinacker2026}, with its moments, and it enters the form at \(n\) of the preceding paragraph as a shift \(B/A\to B/A-q_b^2/q^2\), \(q_b^2=12\hat U\chi^2/m^2(\hat V+2\hat T)\), which is \(\chi^2/m^2\) on \(S^2_N\) at large \(N\). No term \((v\cdot v)(\alpha\cdot\alpha)\) arises, so on the component of \(\alpha\) orthogonal to \(v\) this term gives no mass. Then \(\varrho\to q^2\varrho/(q^2-\varrho q_b^2)\), and at \(V=T\)
\[
-L\varrho_0^2\Bigl[\frac{(\tfrac12\beta_K-g)^2}{q^2-q_b^2+q_*^2}+\frac{g-g_2}{q^2}+\frac{g_2-g^2}{q^2-q_b^2}\Bigr].
\]
So the forms of the preceding paragraph hold for \(q^2\gg|q_b^2|\). Below that scale, for \(\hat U>0\) the static scalar pole has the tachyonic sign unless \(q_*^2\) exceeds \(q_b^2\), and for \(\hat U<0\) and \(\mathcal M\ge0\) it is a Yukawa pole. That work uses \(V_{01}>0\) from \(S_{\mathrm{vac}}=-\int\rho_M\rho^{-2}V_0\) in Equation~(4.46), and in the passages read does not state the sign of \(V_0\); with the same convention the sampled \(K\) with \(\hat U>0\) have \(V_0<0\), normalisations not compared, so no conclusion on that work is drawn. The tree non-derivative terms, the terms that cancel the tadpole of \(U_4\) at \(q=0\), and any term with the same fibre dependence are not computed, and a long-range form needs them or \(\hat U\to0\); where \(\chi^2\lesssim\vartheta\) the second-order non-derivative terms of \(U_2\) and of the first term are of the size of this one and are among those not computed.

\emph{The inputs \(V=T\) and equal weights.} \(V\) and \(T\) both scale as \(\chi^2\), so the response of \(\chi\) does not move the condition \(V=T\). On \(S^2_N\), \((V-T)/T\simeq(4\ln N+1.016)/4N^2\to0\) (numerical), so \(V=T\) is approached at large \(N\) for this \(K\) and does not need a tuned \(K\); it holds exactly at \(f_1=f_2+f_3\). At \(V=T\), \(U_2=-(3F/128\pi^2)\vartheta\hat T\sqrt{|g|}\chi^2<0\), so for \(\hat U\ge0\) and \(Y\le4|U_4|+2|U_2|\), (S1) and \(V=T\) are incompatible at a local minimum: there the sign of the torsion term that supplies the minimum makes \(\beta_K\) non-positive. For \(\hat U\ge0\) and \(Y>4|U_4|+2|U_2|\), \(\beta_K>0\), and its value is set by the tree term of negative slope. On \(S^2_N\) the weights of \(D\) reach \(F\hat V\chi^2\), \(F\hat T\chi^2\) up to a relative correction of order \((n/\chi N)^2\) with a logarithm, with ratios \(0.9994\) at \(N=1024\), \(n=10\), \(\chi=1\), and those of \(B_0,B_1\) carry a common factor from the truncation \(J'\le n\), about \(1-3/2n\) (\(0.8634\) there), the same for \(V_1,T_1,T_2\). So equal weights among the loop levels hold for \(\chi N\gg n\gg1\) on this \(K\). The hypothesis on the external harmonics above \(J=1\) is not tested.

\emph{Consequence.} Within these hypotheses a static, late-time response with ratio \(1\), read in the variable to which the source couples, requires a light \(t\) and \(\beta_K=1\), and the computed exponent is negative for \(\hat U\ge0\) where \(Y<4|U_4|+2|U_2|\). With a light \(t\) and \(\beta_K<0\) the deviation of this ratio from \(1\) exceeds \(1\), and lies in \((1,\tfrac43]\) for \(\beta_K\in[-1,0)\), to be set against the value \((2.1\pm2.3)\times10^{-5}\) of \(\gamma-1\) in~\cite{Bertotti2003}. This is not a Cassini bound on the emergent metric, since solar-system scales lie where \(q^2\ll|q_b^2|\) and the non-derivative terms are not computed. The value \(\beta_K=1\) needs a tree term of negative slope at the minimum and, for \(\hat U\ge0\), \(Y>4|U_4|+2|U_2|\), which at large \(N\) on \(S^2_N\) is \(\chi^2<(\sqrt{39}-3)\vartheta/4\); for \(\hat U<0\), where \(q_b^2<0\), it is available at negligible \(Y\). In each case the relation among the tree data that sets it is not derived, and the vacuum term must be cancelled at the minimum without cancelling its derivative in \(\chi\).

These statements are for the hypotheses of the preceding paragraph, the product background without mixed flux, one loop, the formal continuation and the frozen symbol, for scales of \(q\) as stated; the first-order response in the static channel at \(q\cdot y=0\), with \(V=T\) where the index part alone enters \(U_2\) and with \(W_R\propto\hat V-\hat T\) neglected; the fibre-constant background for \(U_2\) and for the response of \(Y\), the fibre flux in \(Y\) at the background only and with the base-fibre cross term dropped, and higher orders in \(\vartheta/\chi^2\mu_A^2\) not treated; and a mass of the matter modes that comes from \(\chi\) alone, so that masses from gauge dynamics, which see \(g^2e^{-2\phi}\), and the non-singlet moduli of \(K\) are not treated. The sign statement and the exponents are exact for these forms, the first-term identity and its non-negativity are exact on the background, the level form of \(Y\) is the leading form in \(s\) stated there, and the values on explicit \(K\) are exact spectra with numerical fits. They do not construct the \(K\) selected by the theory, determine \(\mathcal M\) or the background value of \(\phi\), compute the loop kinetic term of \(\chi\), the rotating class, higher loops or the fibre Laplacian, identify \(\bar h\), \(h_\beta\) or \(t\) with quantities of the emergent metric of Chapter~\ref{ch:vierbein-emergence}, discharge Assumption~\ref{ass:gap-psi}, or bear on Assumption~\ref{ass:lorentzian}. Problems~C, F and~N remain open and PROGRAMME-CONDITIONAL.

